%% --- Topic 8 Log Sheet --- %%
\chapter*{Topic 8 Log Sheet}

\section*{Project Title}

Music Visualisation Project

\section*{Topic}

Topic 8 - Async and Callbacks

\section*{What progress have you made this topic?}

\subsection*{1. Update the visuals}

\begin{itemize}
  \item Visual 3 - Update the progress bars from five colors to seven colors
  \item Visual 4 - Add two trees which branched varied by the music
  \item Visual 5 - Add the cocktail (with bubbles) and lyrics for the visual
\end{itemize}

\subsection*{2. Preparing the test cases}

\begin{itemize}
  \item set up integration test and unit test cases
\end{itemize}

\subsection*{3. Prepare the report}

\begin{itemize}
  \item study the questions and items to be replied for the report
\end{itemize}

\section*{What problems have you faced and were you able to solve them?}

\subsection*{1. Visualising the music frequency map}

\begin{itemize}
  \item in The Rainbow visual, the fourier.getEnergy function has five string arguments, as per p5.sound documentation
  \item but the rainbow has seven colors
  \item to make the sound frequency spectrum divide by seven spectrum, another argument input method has to be complied
  \item the getEnergy(min, max) input method will generates out energy of sound between min frequency and max frequency
  \item by dividing human-recognisable sound frequencies into seven ranges (ignore 16kHz - 32kHz which is too high for my songs), sound energy can be extracted and being mapped onto the status bar of each of the seven colors
  \item hence the rainbow can be generated
\end{itemize}

\subsection*{2. Adding the trees}

\begin{itemize}
  \item in Garden visual, since the original phyllotaxis part is not filling in the screen
  \item I have added two trees which is out of the mid-term project plan
  \item here, and old programming language Logo is being remake (in the Turtle class)
    \begin{itemize}
      \item a turtle can leave a track (this is drawing) when it is walking on the screen
      \item the turtle can move forward, backward and turn for each step
      \item more detail can make reference to \url{https://en.wikipedia.org/wiki/Logo_(programming_language)}
    \end{itemize}
  \item data structure stack is used to store its steps, for moving backward for drawing branches in other way
  \item hence, this tree, with two branches on each node, can be donw by making the turtle to and fro, with the use of stack pop and stack push functions
  \item below shows recursion codes of the movement of turtle movements
\end{itemize}

\subsection*{3. Adding lyrics}

\begin{itemize}
  \item For making the Disco visual more attaractive, I have added the lyrics show, on the top of the cocktail cup
  \item this is actually calling p5.js loadTable function (this is for data visualisation app), to read in the lyrics and their display time (in ms)
  \item As the song is Korean song, one more font is being used to display the lyrics
\end{itemize}

\section*{Are you on target to successfully complete your project? If you aren’t on target, how will you address the issue?}

\begin{itemize}
  \item My project is a little bit behind the schedule
  \item the terrain in the first visual will be finished in next log report
  \item the test cases and report will start from next log
  \item Async and Await is being delayed due to the making of other visual parts
\end{itemize}

\section*{What are you planning to do over the next few weeks?}

\begin{enumerate}
  \item Visual 1 - North Pole Light
    \begin{itemize}
      \item remove the Mountain visuals
      \item replace by the mountains to the terrain the reflect sound property
    \end{itemize}
  \item Visual 2 - The Sun
    \begin{itemize}
      \item update of this visual has almost been finished
    \end{itemize}
  \item Visual 3 - The Rainbow
    \begin{itemize}
      \item update of this visual has almost been finished
    \end{itemize}
  \item Visual 4 - Garden
    \begin{itemize}
      \item update of this visual has almost been finished
    \end{itemize}
  \item Visual 5 - Disco
    \begin{itemize}
      \item cocktail cup, bubble and lyrics has finished
      \item in next steps I will add the disco lights
    \end{itemize}
  \item Visual 6 - Performance Platform
    \begin{itemize}
      \item after review the project scope and project timeline, this has to be withdrawn from my production plan
    \end{itemize}
\end{enumerate}

Except from buinding up the visuals and the mp3 player, I will also arrange the following

\begin{itemize}
  \item add Async and Await onto visuals
  \item arrange tests to the project
  \item prepare the report
\end{itemize}

\section*{ Appendix - The summary of code I have added}

\subsection*{Make the Rainbow visual more complete}

Below is the capture of The Rainbow that capture sound energy within assigned frequencies spectrum, into energy value for rainbow visual.

\lstinputlisting[language=JavaScript]{./Topic8/script1.js}

\subsection*{Dawing of tree in Garden visual}

Below is the capture of how the tree is being drawn by Logo Turtle.

\lstinputlisting[language=JavaScript]{./Topic8/script2.js}

\subsection*{Showing the Lyrics}

Below is the code for loading csv file to get lyrics and then display them in the Disco visual

\lstinputlisting[language=JavaScript]{./Topic8/script3.js}
